% Original PDF page 17; hash in recovery-map.json.
\recoverypage{17}
\noindent
Disclosure to finalize. LLM assistance was used to research and draft this manuscript and its\linebreak[4]
reference-example code. The author must verify the text, formulas, examples, implementation claims,\linebreak[4]
and bibliography before submission. Record the actual models, dates, and material implementa-\linebreak[4]
tion/formalization/experiment uses when the final runs exist; do not infer a model version or an author\linebreak[4]
review that has not occurred. The reference suite is not a claim of human-independent discovery or\linebreak[4]
audited proof.\par

\noindent
Remaining source and format work. Replace planned-result fields after freezing the source and\linebreak[4]
experimental manifests. Validate actual modal bridges, source-fidelity annotations, positive proof\linebreak[4]
paths, counterexamples, and availability. Prepare the anonymous implementation supplement; keep\linebreak[4]
the private audit with public repository identities outside the submission. Replace the missing official\linebreak[4]
checklist file rather than inventing its questions. The workshop style must remain unchanged; recheck\linebreak[4]
the main-text page count after adding results. All appendices are part of the same submission PDF.\par

\section*{H  Learning-system anatomy and terminology}
\noindent
The original methods draft emphasized generic learned advice. This revision examines the concrete\linebreak[4]
modal autoencoder, packed tensor training, source-free feature promotion, and multi-view targets. All\linebreak[4]
statements here describe the inspected implementation; no new model was trained in this manuscript\linebreak[4]
preparation.\par

% Editable layout transcription; cell boundaries not asserted.
\begin{center}\resizebox{\linewidth}{!}{\begin{tabular}{@{}l@{}}
\hspace*{14\fontdimen2\font}Table 7: Objects that must not all be called the same latent representation. \\
 \\
Object\hspace*{20\fontdimen2\font}Construction and use\hspace*{22\fontdimen2\font}What the current evidence does \\
\hspace*{68\fontdimen2\font}not establish \\
Canonical modal IR\hspace*{8\fontdimen2\font}Source-linked formulas, operators,\hspace*{8\fontdimen2\font}Unsupervised invention of the \\
\hspace*{26\fontdimen2\font}predicates, roles, conditions, excep-\hspace*{5\fontdimen2\font}schema; arbitrary whole-language \\
\hspace*{26\fontdimen2\font}tions and frame structure.\hspace*{16\fontdimen2\font}correctness. \\
Structured round-trip\hspace*{5\fontdimen2\font}Seven-field rule records under a spe-\hspace*{5\fontdimen2\font}Full fidelity for every richer modal \\
IR\hspace*{24\fontdimen2\font}cific compiler, vocabulary and real-\hspace*{6\fontdimen2\font}declaration. \\
\hspace*{26\fontdimen2\font}ization profile. \\
Feature encoding\hspace*{10\fontdimen2\font}Token/cue, semantic-slot, frame, con-\hspace*{5\fontdimen2\font}A pure text-only input when an \\
\hspace*{26\fontdimen2\font}tract, scope, provenance and diagnos-\hspace*{5\fontdimen2\font}existing parse or bridge target is \\
\hspace*{26\fontdimen2\font}tic features.\hspace*{29\fontdimen2\font}used. \\
Decoded embedding\hspace*{9\fontdimen2\font}Numerical prediction compared\hspace*{13\fontdimen2\font}Generated natural-language recon- \\
\hspace*{26\fontdimen2\font}against the sample embedding using\hspace*{8\fontdimen2\font}struction or a guaranteed informa- \\
\hspace*{26\fontdimen2\font}MSE and cosine.\hspace*{27\fontdimen2\font}tion bottleneck. \\
Family/view distribu-\hspace*{5\fontdimen2\font}Class scores over configured modal\hspace*{8\fontdimen2\font}Probability that a statement is \\
tion\hspace*{22\fontdimen2\font}families and compiler/prover view\hspace*{9\fontdimen2\font}legally true or a theorem is proved. \\
\hspace*{26\fontdimen2\font}components. \\
Source-withheld text\hspace*{6\fontdimen2\font}Deterministic realization from admit-\hspace*{5\fontdimen2\font}Proof of original-source fidelity; \\
\hspace*{26\fontdimen2\font}ted IR without originating surface ac-\hspace*{4\fontdimen2\font}compiler and decoder can share \\
\hspace*{26\fontdimen2\font}cess.\hspace*{37\fontdimen2\font}omissions. \\
Proof-feedback heads\hspace*{6\fontdimen2\font}Bounded supervised predictions from\hspace*{7\fontdimen2\font}A substitute for the next actual \\
\hspace*{26\fontdimen2\font}admitted, version-bound verifier feed-\hspace*{4\fontdimen2\font}proof check. \\
\hspace*{26\fontdimen2\font}back. \\
Stable guidance ex-\hspace*{7\fontdimen2\font}Selected shared features mapped to\hspace*{8\fontdimen2\font}Raw parameter state becoming a \\
port\hspace*{22\fontdimen2\font}canonical view contracts after canary\hspace*{5\fontdimen2\font}new semantic authority. \\
\hspace*{26\fontdimen2\font}admission.
\end{tabular}}\end{center}

\noindent
A useful classification. The concrete adaptive component is a feature-based, multi-head recon-\linebreak[4]
struction/classification learner with sparse state and a packed tensor execution path. It should not\linebreak[4]
be described as a variational autoencoder, generative adversarial network, Transformer sequence-\linebreak[4]
to-sequence formalizer, or trained universal semantic bottleneck without identifying a separate\linebreak[4]
implemented model and checkpoint. Other experimental mechanisms may exist elsewhere; this\linebreak[4]
classification applies to the exact paths audited here.\par

